% Pursuing Stacks (cotes 134-2 to 134-8): archive cote 134-6, site page 53 (PDF page 54), typescript page 371.
% typescript. Draft compared: ps5.tex:4004-4060.
% AI-transcribed by Claude Opus 5.5 (claude-opus-5-5, adaptive thinking, effort high), Message Batches API, batch ps3 (msgbatch_01XdFmuVVBAcgc7EFxdtCqsj), 2026-11-12; not yet checked by eye.
\scanpage{134-6/53}
\mspage{371}

know -  although a fair number of \inserted[typed]{partly} overlapping results in this direction
are known,  among the oldest being the relevant ``universal coefficients
formulae''  relating homology and cohomology (reducing all to a formula of
the type (2) or (10) above),  or Cartan's old seminar on \xout{sheaf} Leray's
sheaf theory,  \xout{stating that} introducing singular homology with coefficients
in a sheaf and proving that on a topological variety,  this was (up to
dimension shift and twist by the \xout{orientation} twisted integers) essentially
the same as singular cohomology\inserted{(with coefficients in sheaves too).} It is not sure that an all-inclusive sta-%
tement of equivalence between homology and cohomology (in those situations
when such equivalence is felt to hold indeed) does at all exist  - \xout{it may}
at any rate,  according to what kind of coefficients \xout{are} one wants to con-%
sider,  and what kind of \inserted[typed]{extra} structures one is interested in when dealing with
homology and cohomology invariants,  it seems that each of the two points
of view has an originality \inserted[typed]{and advantages,} of it's own and cannot be entirely \xout{subsumed} \inserted[typed]{superseded} by
the other.  From the contexts I have been mainly working in,  there defi-%
nitely was no choice,  namely cohomology (including non commutative one)
was the basic data,  while sheaves and their generalizations (such as com-%
plexes of sheaves,  or stacks)  w\over{\unsure{he}}{e}re the coefficients.  I do'nt remember
of any moment where I would have paused and asked myself \uline{why} in most con-%
texts where I was working in (whose common denominator was topoi),  there
was'nt any direct hold on anything like homology invariants.  The reason
for this, \inserted{inertness of mine,} probably,  is that the cohomology formalisms I hit upon were
self-contained enough,  so as to leave no regret for the absence of a
homology formalism,  or at any rate of a more or less direct description
of it independently of cohomology.  Another reason,  surely,  is that I
did'nt have too much contact with topologists and homotopists and their
everyday \xout{work} tools,  such as Steenrod operations,  homology of the symme-%
tric groups,  and the like.  This question of ``why this reluctance of
homology to show forth''  has finally \xout{come} surfaced only during these very
last days,  when the answer for it (or one possible answer at any rate) is
becoming evident:  namely,  that \uline{for a general topos},  embodied by a
category of sheaves (of sets) \hand{$\mathcal{A}$} ,  \uline{there are not enough projectives in}
\uline{\hand{$\mathcal{A}$} ,  and not even enough projectives in \hand{$\mathcal{A}$}$_{\mathrm{ab}}$} ,  the category of abelian
sheaves\unsure{.}  It is becoming appearant (what surely everybody has known for
ages)  that \inserted[typed]{in technical terms,} \uline{doing ``homology''} \uline{is working with p\over{\illegible{1}}{ro}jectives},  \uline{while doing}
\uline{``cohomology'' is working with injectives.}  As the\over{\unsure{fe}}{re} are enough injectives
in \hand{$\mathcal{A}$}$_{\mathrm{ab}}$  but not enough projectives,  cohomology is around and homology
not,  period !

There is however a rather interesting class of topoi admitting suff\unsure{i}-%

\margin{53}
